\documentclass[12pt,a4paper]{article}

\usepackage{polyglossia}
\setmainlanguage{french}

\usepackage{amsmath, amssymb}

\usepackage{graphicx}
\usepackage{caption}

\usepackage{siunitx} % Pour les unités

\usepackage{geometry}
\geometry{top=2cm, bottom=2.5cm, left=3cm, right=3cm}
\setlength{\parindent}{0pt}

\usepackage{enumitem}

\usepackage{hyperref}

\usepackage{xcolor}

\usepackage{multicol}
\setlength{\columnseprule}{1pt}
\def\columnseprulecolor{\color{gray}}

\usepackage{tcolorbox}


\usepackage{import}
\usepackage{xifthen}
\usepackage{pdfpages}
\usepackage{transparent}
\usepackage{float}

\newcommand{\incfig}[1]{%
	\def\svgwidth{\columnwidth}
	\import{./figures/}{#1.pdf_tex}
}


\newcommand{\RomanNum}[1] %pour les chiffre romain
{\MakeUppercase{\romannumeral #1}}

\usepackage{myunit}
\usepackage{stanli}








\title{Mur de Soutènement}
\author{}
\date{\today}

\begin{document}
	\counterwithin{section}{part} % pour que \section{} recommence de 1 aprés \part{}
	
	
	\tableofcontents
	
	\newpage
	
	
	
	\part{Caractéristiques fondamentale du sol}
	
	
	\section{Densité $\gamma$}
	
	Elle conditionne la valeur de la contrainte verticale $\sigma_{v}$ à un profondeur Z de la surface du sol :
	
	\[ \sigma_{v} = \gamma \cdot Z \]
	
	Elle souvent prise entre $18\knmc$ et $20\knmc$
	
	
	\section{Portance du sol q}
	
	C'est la contrainte de calcul que peut reprendre le sol sans se déformer.
	
	\subsection{Valeurs de la portance q}
	
	\begin{table}[H]
		\begin{center}
			\begin{tabular}{l|c}
				Nature du sol & $q\mpa$ \\ \hline
				Roches saine (non fissurée) & 0,75 à 4,5 \\
				Sable compact & 0,35 à 0,75 \\
				sable lâche (peu compact) & 0,2 à 0,4 \\
				argile & 0,1 à 0,3
			\end{tabular}
		\end{center}
	\end{table}
	
	
	
	\section{Angle et l cohésion du sol}
	
	\subsection{Angle de talus naturel $\varphi$}
	
	C'est l'angle maximal (avec l'horizontal) que peut présenter un sol mis simplement en tas.
	
	\subsection{Cohésion C}
	
	C'est la capacité qu'un sol s'amalgame (coller).
	
	Sol pulvérulent $\rightarrow$ C = 0
	
	\subsection{Valeurs de $\varphi$ et de C}
	
	
	
	\begin{table}[H]
		\begin{center}
			\begin{tabular}{l|c|c}
				Nature du sol & $C ~\unit{\kilo\pascal}$ & $\varphi\degree$ \\ \hline
				Gravier & 0 & 40 à 50 \\
				Sable compact & 0 à 10 & 30 à 40 \\
				Sable lâche & 0 à 10 & 25 à 30 \\
				Argile & 20 & 15 a 25
			\end{tabular}
		\end{center}
	\end{table}
	
	
	
	
	
	\part{Pré-dimensionnement}
	
	
	
	\begin{center}
		\begin{tikzpicture}[scale=.9]
			
			
			\coordinate (a) at (0,0);
			\coordinate (b) at (13,0);
			\coordinate (c) at (13,2);
			\coordinate (d) at (9,2);
			\coordinate (e) at (9,9);
			\coordinate (f) at (6,9);
			\coordinate (g) at (4,2);
			\coordinate (h) at (0,2);
			
			\coordinate (i) at (-2,5);
			\coordinate (j) at (5,5);
			\coordinate (k) at (15,9);
			
			\coordinate (m) at (4,0);
			\coordinate (n) at (9,0);
			
			
			\path (f) -- (g) coordinate[pos=0.3] (l);
			
			\draw[thick] (a) -- (b) -- (c) -- (d) -- (e)
			-- (f) -- (g) -- (h) -- (a);
			
			\draw[dashed] (i) -- (j) node[above,midway]{T.N};
			\draw[dashed] (e) -- (k) node[above,midway]{T.N};
			
			
			\draw[thick,<-] (l) -- ++(150:2) node[above,left]{Fruit $\min 2 ~\unit{\percent}$};
			
			
			
			
			\dimensioning{1}{f}{e}{10}[$\dfrac{H}{24}$];
			\dimensioning{1}{a}{b}{-1}[$ 0,2H \leq B \leq 0,6H$];
			\dimensioning{1}{a}{m}{-2.3}[$b_{2} = \dfrac{B}{3}??$];
			\dimensioning{1}{m}{n}{-2.3}[$\dfrac{H}{12}$];
			\dimensioning{1}{n}{b}{-2.3}[$b_{1}$];
			
			\dimensioning{2}{i}{a}{-1}[$h_{2}$];
			\dimensioning{2}{e}{c}{14}[$h_{1}$];
			\dimensioning{2}{c}{b}{14}[$\dfrac{H}{12}$];
			
			\dimensioning{2}{k}{b}{15.3}[$H$];
			
			
			
			
		\end{tikzpicture}
	\end{center}
	
	
	
	
	\part{Calculs des forces horizontales}
	
	
	\section{Poussée}
	
	Contrainte horizontale : 
	
	\[ \sigma_{h} = K_{a} \cdot \gamma \cdot Z \]
	
	Avec $K_{a} = \tan^{2}(\dfrac{\pi}{4} - \dfrac{\varphi}{2})$ Coefficient de poussée.
	
	
	\begin{multicols}{2}
		
			L'effort résultant : 
			
			\[ P = \dfrac{K_{a} \cdot \gamma \cdot H \cdot H}{2} \]
			
			Cette effort agit a $\dfrac{H}{3}$
			
			\columnbreak
			
			\begin{center}
				\begin{tikzpicture}[scale=.3]
					
					
					
					\coordinate (a) at (0,0);
					\coordinate (b) at (13,0);
					\coordinate (c) at (13,2);
					\coordinate (d) at (9,2);
					\coordinate (e) at (9,9);
					\coordinate (f) at (6,9);
					\coordinate (g) at (4,2);
					\coordinate (h) at (0,2);
					
					\coordinate (i) at (-2,5);
					\coordinate (j) at (5,5);
					\coordinate (k) at (15,9);
					
					\coordinate (m) at (4,0);
					\coordinate (n) at (9,0);
					
					
					\path (f) -- (g) coordinate[pos=0.3] (l);
					
					\draw[thick] (a) -- (b) -- (c) -- (d) -- (e)
					-- (f) -- (g) -- (h) -- (a);
					
					\lineload{1}{e}{d}[0][5];
					
					
				\end{tikzpicture}
			\end{center}
	
	\end{multicols}
	
	
	\section{Butée}
	
	Contrainte horizontale : 
	
	\[ \sigma_{h} = K_{p} \cdot \gamma \cdot Z \]
	
	
	avec $K_{p} = \tan^{2}(\dfrac{\pi}{4} + \dfrac{\varphi}{2})$ étant le coefficient de butée.
	
	\begin{multicols}{2}
		
	L'effort résultant : 
	
	$b = \dfrac{h \cdot K_{p} \cdot \gamma \cdot h}{2}$
	
	Situé a $\dfrac{1}{3}h$
	
	\columnbreak
	
	
	\begin{center}
		\begin{tikzpicture}[scale=.3]
			
			\coordinate (a) at (0,0);
			\coordinate (b) at (13,0);
			\coordinate (c) at (13,2);
			\coordinate (d) at (9,2);
			\coordinate (e) at (9,9);
			\coordinate (f) at (6,9);
			\coordinate (g) at (4,2);
			\coordinate (h) at (0,2);
			
			\coordinate (i) at (-2,5);
			\coordinate (j) at (5,5);
			\coordinate (k) at (15,9);
			
			\coordinate (m) at (4,0);
			\coordinate (n) at (9,0);
			
			
			\path (f) -- (g) coordinate[pos=0.3] (l);
			
			\draw[thick] (a) -- (b) -- (c) -- (d) -- (e)
			-- (f) -- (g) -- (h) -- (a);
			
			\coordinate (A) at (5,0);
			\lineload{1}{j}{A}[0][-6];
		
			
		\end{tikzpicture}
	\end{center}
	
	\end{multicols}
	
	
	\section{Surcharge sur le terrain}
	
	
	Contrainte horizontale :
	
	\[ \sigma_{h} = (K_{a} \text{ ou } K_{p}) \cdot q \]
	
	L'effort résultant : 
	
	\[ Q = (K_{a} \text{ ou } K_{p}) \cdot q \cdot H \]
	
	
	\begin{tcolorbox}[colback=gray!20, colframe=black, title=Remarque]
		En générale, on ne tient pas compte par mesure de sécurité de la butée B, qui peut être mobilisée qu'après un certain déplacement du mur (effort difficilement mobilisable)
	\end{tcolorbox}
	
	
	
	
	
	
	
	
	
	
	
	
	
	
	
	
	
	
	
	
	
	
	
	
\end{document}